\begin{houseviewbox}[结论先行]
本报告不接受“资金已经流入”“仓位太高所以必须反弹”这类不可复算前提。产业层面，半导体设备、海外 BD 和电网建设分别有真实的资本开支、合同或政策资料；证券层面，六个研究样本均未形成客户--订单--利润--现金的闭环。故现阶段最优行动不是横向追主题，而是维护一个有明确升级条件的观察池。
\end{houseviewbox}

\begin{exhibitbox}[主题结论与行为边界]
\small
\begin{tabularx}{\exhibitboxwidth}{L{2.3cm}X L{2.1cm}X}
\toprule
主题 & 已验证的事实 & 当前标签 & 触发升级的最低条件\\ \midrule
半导体设备 & 海外设备商披露先进逻辑、存储升级和先进封装带动投资及 backlog。 & 市场支持观察 & 中国样本的产品、客户认证、订单/验收、ASP、毛利和现金回款。\\
创新药 BD & 恒瑞--GSK 和信达披露证明首付款、里程碑、合同负债与收入确认可显著错期。 & 选择性跟踪 & 单资产会计、成功概率、后续成本、销售提成和每股收益桥。\\
电力设备 & 能源主管部门披露电网/新能源建设与新型电力系统试点。 & 宏观支持观察 & 中标份额、在手订单、交付验收、价格/成本、应收和经营现金流。\\
\bottomrule
\end{tabularx}
\end{exhibitbox}
\sourcenote{SRC-05--08、SRC-09--10；行业与供应链分析。}

\section{研究样本与当前行动}
半导体设备观察北方华创（002371）与中微公司（688012）；创新药观察恒瑞医药（600276）与百利天恒（688506）；电力设备观察思源电气（002028）与国电南瑞（600406）。这六家公司是为测试文章逻辑而设的样本，不是推荐名单。全链条审计将全部列为 \stance{riskred}{watchlist only / insufficient evidence}，原因不是否定其商业价值，而是本案例尚未取得能支持当前价格估值的完整公司包。

\begin{riskbox}[本报告不能回答的问题]
截至资料截止日，缺少可复现的 2026-07-24 样本价格、流动性及资金流表，亦缺少六家公司可核验的财务、客户和订单包。因此不存在可信的目标价、上涨空间或仓位建议；用单日资金流、行业报告标题或潜在 BD 总金额填补这些缺口会造成虚假精度。
\end{riskbox}
